\section*{Overview}

Aho-Corasick is the multi-pattern string matcher for ``one text, many patterns.'' It starts from a trie of all
patterns, then adds failure links so the scan never restarts from scratch after a mismatch.

\section*{When to Use It}

Use it when:

\begin{itemize}
  \item there are many patterns,
  \item the text is long enough that running KMP or naive search for each pattern separately is too slow,
  \item shared pattern prefixes are significant.
\end{itemize}

\section*{Core Idea}

Build a trie of the patterns. For each trie node, compute a failure link to the longest proper suffix of that node's
string that is also a trie prefix. During text scanning:

\begin{itemize}
  \item follow the next trie edge when it exists,
  \item otherwise follow failure links until a valid transition exists,
  \item report all pattern endings attached to the reached state.
\end{itemize}

\includegraphics[width=\textwidth]{figures/failure-links.svg}

\section*{Key Insight}

The failure link is the multi-pattern version of the KMP fallback. It preserves the longest suffix of the processed text
that could still grow into some pattern. That is why already-checked characters are never rescanned.

\section*{Worked Problem}

\paragraph{Problem.}
You are given a dictionary of words and one long text. Count how many times each dictionary word appears in the text.

\paragraph{Why Aho-Corasick fits.}
All patterns are queried against the same text. Their shared prefixes belong in one trie, and failure links prevent
restarting from the root after every mismatch.

\paragraph{Scan view.}
If the current state spells \texttt{hers} and the next character breaks that path, the automaton jumps by failure links
to the longest suffix that is still a valid trie prefix, for example \texttt{s} or \texttt{he}, depending on the
dictionary.

\section*{Correctness Intuition}

At every position, the automaton state represents the longest suffix of the processed text that is also a trie prefix.
Normal transitions extend that suffix; failure links shorten it to the next-best valid candidate. Therefore every
pattern ending at the current position is exposed either directly in this state or through output information inherited
from failure ancestors.

\section*{Complexity Analysis}

If total pattern length is \(m\) and text length is \(n\), construction plus scanning is linear in \(m+n\) up to the
chosen transition representation.

\section*{Implementation}

The code uses a lowercase alphabet trie and supports:

\begin{itemize}
  \item pattern insertion,
  \item failure-link construction,
  \item text scanning that counts occurrences of every pattern.
\end{itemize}

\section*{Common Pitfalls}

\begin{itemize}
  \item Forgetting to propagate output information along failure links.
  \item Using a dense transition table on a huge alphabet and wasting memory.
  \item Treating the automaton state as if it stored only one matched pattern.
\end{itemize}

\section*{Variants / Extensions}

\begin{itemize}
  \item DP on the automaton for forbidden-substring problems.
  \item Sparse transition maps for larger alphabets.
  \item Counting total matches, distinct matches, or first/last occurrence positions.
\end{itemize}

\section*{Practice Problems}

\begin{itemize}
  \item Count how many times each pattern appears in a long text.
  \item Mark all positions where any forbidden pattern occurs.
  \item DP on strings that must avoid a set of bad substrings.
\end{itemize}

\section*{References}

\begin{itemize}
  \item \href{https://cp-algorithms.com/string/aho_corasick.html}{cp-algorithms: Aho-Corasick}.
  \item \href{https://usaco.guide/adv/string-search}{USACO Guide: String Searching}.
  \item \href{https://codeforces.com/catalog}{Codeforces Catalog}.
\end{itemize}
